
\exercise{An unusual case\label{exBrainUnusual}}{4}
Determine the type of the bifurcation that occurs in the system
\eqn{
\dot{x}=(x-1-p)(2p-x)
}
Check your answer carefully.

\solution 
First we need to find the bifurcation. In this case this is very easy as the equation of motion is already factorized. Hence the steady states are 
\eqa{
x_1^* &=& 1+p \\
x_2^* &=& 2p
}
Because we don't see any square roots the bifurcation probably won't be fold. Also since we only have one variable it can't be a Hopf. Moreover considering the steady states we see that the two branches cross in $p=1$, $x^=2$. This sort of crossing looks very much like a transcritical bifurcation, except for one small issue: neither of the two steady states remains pinned at a particular value. 

In this chapter we have shown that pinning of a steady state to a particular value leads to transcritical bifurcations, but is the reverse true as well? Does the transcritical bifurcation actually require the pinning. 

Let's check this carefully. the right side of our differential equation is 
\eq{
f(x,p) = (x-1-p)(2p-x)
}
We claim that at $p=1$, $x^*=2$ is a steady state. Indeed, 
\eq{
f(2,1) = (2-1-1)(2\cdot 1 -2) = (0)(0) = 0. \quad \mbox{(Stationarity Condition)}
}
Moreover we check
\eqa{
f_{\rm x}(2,1) &=& \left. \partial_x (x-1-p)(2p-x) \right|_{x=2,p=1} \\
  &=& \left. (2p-x)-(x-1-p) \right|_{x=2,p=1} \\
  &=& 0-0 = 0
}
which confirms that $p=1$ is a bifurcation. Normally we would expect this to indicate a fold bifurcation, so let's show that one of the generaicity conditions of the fold is broken. We compute
\eqa{
f_{\rm p} &=&   \left. \partial_p (x-1-p)(2p-x) \right|_{x=2,p=1} \\ 
  &=&   \left. -(2p-x)+2(x-1-p) \right|_{x=2,p=1} \\
  &=&  -0+2(0) = 0
}
Hence the genericity condition of the fold bifurcation $f_{\rm p}\neq 0$ is broken. As we already suspected, this proves that this in not the fold. The two conditions 
\eq{
   f_{\rm p}=0 \qqq f_{\rm x}=0 
}
are the bifurcation conditions of the transcritical bifurcation, but is it really a transcritical or something more special? We can find out by testing the genericity conditions of the transcritical bifurcation, $f_{\rm xp}\neq 0$ and $f_{xx}\neq 0$. 

We already computed 
\eq{
\partial_xf = (2p-x)-(x-1-p)
}
which we can now use to compute 
\eqa{
f_{\rm xx} &=& \left. \partial_x (2p-x)-(x-1-p) \right|_{x=2,p=1} \\
  &=& -2 \neq 0
}
and also 
\eqa{
f_{\rm xp} &=& \left. \partial_p (2p-x)-(x-1-p) \right|_{x=2,p=1} \\
  &=& 2-1 = 1 \neq 0
}
as the two genericity conditions are met, this really is a transcritical bifurcation, albeit an unusual one. 
